\documentclass[article]{jss}

%% -- LaTeX packages and custom commands ----------------------------------------
\usepackage{amsmath,amssymb,amsthm}
\usepackage{algorithm}
\usepackage{algpseudocode}
\usepackage{booktabs}
\usepackage{graphicx}
\usepackage{hyperref}
\usepackage{natbib}
\usepackage{xcolor}

%% -- Theorem environments ------------------------------------------------------
\newtheorem{definition}{Definition}
\newtheorem{proposition}{Proposition}
\newtheorem{remark}{Remark}

%% -- JSS author information ----------------------------------------------------
\author{Danilo A. Sarti\\University / Institution}
\Plainauthor{Sarti}
\title{\pkg{lepidoptera}: A Domain-Agnostic Engine for Causal Simulation with Mixed Variable Types, Vine Copula Dependence, and Counterfactual Propagation}
\Plaintitle{lepidoptera: Causal Simulation with Mixed Types and Vine Copulas}
\Shorttitle{\pkg{lepidoptera}: Causal Simulation Engine}

\Abstract{
We introduce \pkg{lepidoptera}, an open-source R and Python package that provides a modular, domain-agnostic engine for structural causal model (SCM) simulation with mixed variable types. Unlike existing tools that restrict mechanisms to continuous variables or treat interventions as post-hoc operations, \pkg{lepidoptera} elevates the mechanism --- the generating rule associated with each node in a directed acyclic graph (DAG) --- to a first-class, inspectable, and serializable object. The framework natively supports nine type-pair combinations (continuous, categorical, and ordinal in any parent-child configuration), an extensible library of pre-built mechanisms including multinomial logit, ordered logit, group-effect, and categorical transition models, and a vine copula mechanism that learns arbitrary dependence structures from observed data via the Rosenblatt transform. A key design feature is \emph{partial propagation}: by freezing the random draws of unaffected nodes and replaying only the downstream consequences of an intervention, \pkg{lepidoptera} enables individual-level counterfactual reasoning without additional assumptions beyond those of the original SCM. An AI-native extension layer allows any pre-fitted predictive model --- random forests, gradient boosted trees, neural networks --- to serve as a generative mechanism through a single \code{predict\_fn} interface. Comprehensive simulation studies validate type-combination coverage, copula family fidelity, and counterfactual consistency. Real-data illustrations span clinical trial planning, hospital emergency-ward digital twinning, and epidemiological mediation analysis. \pkg{lepidoptera} is available from CRAN and PyPI under the MIT licence.
}

\Keywords{structural causal model, directed acyclic graph, mixed variable types, vine copula, counterfactual simulation, partial propagation, R, Python}
\Plainkeywords{structural causal model, directed acyclic graph, mixed variable types, vine copula, counterfactual, R, Python}

\Address{
  Danilo A. Sarti\\
  E-mail: \email{daniloasarti@gmail.com}
}

%% -- Main document -------------------------------------------------------------
\begin{document}

\section{Introduction}
\label{sec:intro}

Simulation from structural causal models (SCMs) has become an essential tool across empirical disciplines. Epidemiologists use it to plan randomised trials under realistic confounding \citep{hernan2020causal}; economists employ it to evaluate counterfactual policy scenarios \citep{angrist2009mostly}; clinical data scientists use it to stress-test machine-learning pipelines before deployment \citep{vanderschueren2023}; and computer scientists increasingly rely on synthetic data with known causal structure to benchmark causal discovery algorithms \citep{spirtes2000causation, glymour2019review}. Despite this broad demand, the software landscape is fragmented: tools either (i) restrict generating mechanisms to continuous Gaussian variables \citep{scutari2010bnlearn}, (ii) focus on DAG visualisation without generative capability \citep{textor2016dagitty}, (iii) hard-wire longitudinal study designs \citep{sofrygin2017simcausal}, or (iv) conflate the causal structure with a specific model family such as structural equation models \citep{holst2013lava}. No existing package provides a unified, domain-agnostic engine in which the mechanism at every node is an explicit, replaceable object, interventions are first-class operations on that object graph, mixed variable types are supported without workarounds, and individual-level (counterfactual) and population-level (interventional) queries are cleanly separated.

We present \pkg{lepidoptera}, an open-source package available in R and Python with a shared conceptual API, that fills this gap. The name is a nod to the systematic, modular beauty of biological taxonomy: just as a lepidopterist classifies organisms by explicit, composable features, the package classifies data-generating processes by their constituent mechanisms, each typed and independently swappable.

The central contribution is architectural: \pkg{lepidoptera} separates \emph{structure} (which variables exist and who depends on whom, encoded in a DAG) from \emph{mechanism} (how each variable is generated given its parents, encoded in an explicit \code{Mechanism} object). This separation, inspired by the modularity property of SCMs \citep{pearl2009causality, peters2017elements}, has three practical consequences. First, any node mechanism can be replaced for any intervention or sensitivity analysis without re-specifying the rest of the system. Second, partial propagation --- recomputing only the downstream subgraph of an intervention while preserving the random draws of all unaffected nodes --- is exact and requires no additional assumptions, yielding individual-level counterfactual queries \citep{pearl2009causality, shpitser2007identification}. Third, the mechanism slot can hold any callable that respects a minimal signature, so pre-fitted machine-learning models plug in directly via \code{model\_mechanism()}.

The paper is organised as follows. Section~\ref{sec:related} reviews related work. Section~\ref{sec:framework} gives formal definitions. Section~\ref{sec:mechanisms} describes the mechanism library. Section~\ref{sec:simulation} covers the simulation and propagation algorithms. Section~\ref{sec:counterfactual} formalises counterfactual propagation. Section~\ref{sec:copula} presents the vine copula mechanism. Section~\ref{sec:ai} describes the AI-native extension layer. Section~\ref{sec:sims} reports simulation studies. Section~\ref{sec:applications} presents real-data applications. Section~\ref{sec:discussion} concludes.

\section{Related Work}
\label{sec:related}

\subsection{Causal simulation tools}

\citet{sofrygin2017simcausal} introduced \pkg{simcausal} for R, which provides a rich language for specifying longitudinal observational and randomised-study designs with time-varying treatments and censoring. The package has been widely used in causal inference methodology papers \citep{luedtke2016, laan2018} and represents the closest predecessor to \pkg{lepidoptera}. However, \pkg{simcausal} is designed around a temporal index: nodes are time-stamped and generating formulae are specified as R expressions evaluated in a shared environment, which makes it natural for longitudinal settings but unwieldy for cross-sectional DAGs with mixed variable types and heterogeneous mechanisms. Crucially, there is no concept of a reusable, typed mechanism object, and counterfactual (same-individual) queries are not natively supported.

\citet{goldfeld2020simstudy} provide \pkg{simstudy} for R, an expressive package for simulating study populations under a variety of distributional assumptions. Its formula-based interface is convenient for power calculations and trial design, but the package does not expose a causal graph layer, so intervention propagation must be managed manually by the user.

\pkg{bnlearn} \citep{scutari2010bnlearn} focuses on Bayesian network structure learning and inference. It supports simulation from fitted network objects, but mechanisms are restricted to linear Gaussian or discrete conditional probability tables; there is no support for ordinal variables with proportional-odds structure, vine copula dependence, or partial propagation.

For Python, \pkg{causalnex} \citep{beaumont2021causalnex} and \pkg{dowhy} \citep{sharma2020dowhy} provide causal inference pipelines, but neither is designed as a simulation engine: \pkg{causalnex} focuses on structure learning and Bayesian network inference, while \pkg{dowhy} focuses on identification and estimation from observational data, not on data generation.

\pkg{CausalPy} \citep{martinezrodriguez2023} and \pkg{pgmpy} \citep{ankan2015pgmpy} provide Bayesian network tools in Python but share the same limitations: Gaussian or discrete-only mechanisms, no ordinal type, no vine copula, and no partial propagation.

\subsection{Structural causal models and counterfactuals}

The formal foundations of our framework draw on Pearl's structural causal model (SCM) theory \citep{pearl2009causality}, in which a system is defined by a set of structural equations $X_i = f_i(\mathbf{Pa}_i, U_i)$ with exogenous noise terms $U_i$. The \emph{do-operator} $P(Y \mid do(X=x))$ captures population-level interventions by replacing the structural equation for $X$ with a constant \citep{pearl1995causal}. Individual-level counterfactuals $Y_{X=x}(u)$ --- the value $Y$ would have taken for unit $u$ under intervention $do(X=x)$ --- require the additional step of ``abducing'' the noise realisation $u$ from observed data, then applying the intervention in the modified system \citep{pearl2009causality}. \pkg{lepidoptera} operationalises this three-step procedure (abduction, action, prediction) through partial propagation: because the RNG stream for each unit is seeded deterministically from the unit index, the noise realisation is implicitly preserved across baseline and counterfactual runs.

The potential outcomes framework of \citet{rubin1974estimating} and \citet{neyman1923} provides complementary formalisation. The stable unit treatment value assumption (SUTVA) \citep{rubin1980randomization}, which rules out interference and hidden versions of treatment, is satisfied by construction in \pkg{lepidoptera}'s closed-world model.

\subsection{Mixed variable types in graphical models}

Mixed graphical models combining continuous and discrete variables have a long history \citep{lauritzen1989graphical, cox1993linear}. In the context of directed models, \citet{heckerman1995learning} and \citet{lauritzen2002chain} discuss score-based and constraint-based learning for mixed DAGs. However, the focus of that literature is on inference and structure learning, not on user-facing simulation. The practical problem of generating mixed-type samples from a user-specified DAG --- which requires a type-dispatch system mapping every parent-type/child-type pair to a sensible default mechanism --- has received little software attention.

\citet{nagler2016evading} and \citet{czado2019analyzing} provide the theoretical machinery for vine copula construction that underpins \pkg{lepidoptera}'s copula mechanism. The \pkg{rvinecopulib} package \citep{nagler2021rvinecopulib} implements these methods in a high-performance C++ library, which we wrap rather than re-implement.

\subsection{Software infrastructure for reproducible causal analysis}

The \pkg{dagitty} package \citep{textor2016dagitty} provides a web and R interface for drawing and analysing DAGs (identifying adjustment sets, testing d-separation), but it does not generate data. \pkg{lavaan} \citep{rosseel2012lavaan} fits and simulates from structural equation models but is restricted to linear Gaussian mechanisms. \pkg{OpenMx} \citep{neale2016openmx} extends this to more general SEM structures but remains in the continuous-Gaussian paradigm.

Table~\ref{tab:comparison} summarises the feature matrix of the closest tools.

\begin{table}[ht]
\centering
\caption{Feature comparison of causal simulation tools. ``Mixed types'' = native support for continuous, categorical, and ordinal variables in any combination. ``Copula'' = vine copula mechanism. ``Partial prop.'' = individual-level counterfactual via partial propagation. ``R/Py'' = availability in R and Python.}
\label{tab:comparison}
\small
\begin{tabular}{lcccccc}
\toprule
Package & Mixed types & Copula & Partial prop. & Typed mech. & R/Py \\
\midrule
\pkg{simcausal} \citep{sofrygin2017simcausal}  & No  & No  & No  & No  & R \\
\pkg{simstudy} \citep{goldfeld2020simstudy}    & Partial & No & No & No & R \\
\pkg{bnlearn} \citep{scutari2010bnlearn}       & Partial & No  & No  & No  & R \\
\pkg{dagitty} \citep{textor2016dagitty}         & No (viz only) & No & No & No & R \\
\pkg{lavaan} \citep{rosseel2012lavaan}          & No  & No  & No  & No  & R \\
\pkg{causalnex} \citep{beaumont2021causalnex}  & Partial & No & No & No & Py \\
\pkg{dowhy} \citep{sharma2020dowhy}            & No (estim.) & No & No & No & Py \\
\pkg{lepidoptera} (this paper)                 & Yes & Yes & Yes & Yes & R+Py \\
\bottomrule
\end{tabular}
\end{table}

\section{The \pkg{lepidoptera} Framework}
\label{sec:framework}

\subsection{Formal definitions}

Let $\mathcal{V} = \{V_1, \ldots, V_p\}$ be a finite set of variables. Each variable has a \emph{type} $\tau(V_i) \in \{\text{continuous}, \text{categorical}, \text{ordinal}\}$ and a \emph{domain} $\mathcal{D}(V_i)$: an interval $[a, b] \subseteq \overline{\mathbb{R}}$ for continuous variables, a finite unordered set $\mathcal{L}$ for categorical variables, and a finite ordered set $\mathcal{L}$ with a total order $\prec$ for ordinal variables.

\begin{definition}[Mechanism]
A \emph{mechanism} for variable $V_i$ with parent set $\mathbf{Pa}_i \subseteq \mathcal{V} \setminus \{V_i\}$ is a triple $\mathcal{M}_i = (\text{fn}_i, \boldsymbol{\tau}_{\mathbf{Pa}}, \tau_i)$ where:
\begin{enumerate}
  \item $\text{fn}_i: \prod_{V_j \in \mathbf{Pa}_i} \mathcal{D}(V_j) \times \Omega \to \mathcal{D}(V_i)$ is a measurable generating function, with $\Omega$ a probability space carrying all exogenous randomness.
  \item $\boldsymbol{\tau}_{\mathbf{Pa}} = (\tau(V_j))_{V_j \in \mathbf{Pa}_i}$ is the declared parent type vector, used for validation.
  \item $\tau_i = \tau(V_i)$ is the declared output type.
\end{enumerate}
For root variables ($\mathbf{Pa}_i = \emptyset$), $\text{fn}_i$ takes no parent arguments and is called a \emph{marginal mechanism}.
\end{definition}

\begin{definition}[System]
A \emph{system} $\mathcal{S} = (\mathcal{V}, \mathcal{E}, \{\mathcal{M}_i\}_{i=1}^p)$ consists of:
\begin{enumerate}
  \item A finite set of typed features $\mathcal{V}$.
  \item A set of directed edges $\mathcal{E} \subseteq \mathcal{V} \times \mathcal{V}$ forming a DAG $G = (\mathcal{V}, \mathcal{E})$.
  \item A mechanism $\mathcal{M}_i$ for each $V_i \in \mathcal{V}$, such that the declared parent types of $\mathcal{M}_i$ are consistent with the types of $\{V_j : (V_j, V_i) \in \mathcal{E}\}$.
\end{enumerate}
A system is \emph{valid} if the DAG is acyclic, every variable has exactly one mechanism, and the type consistency condition holds.
\end{definition}

\begin{definition}[Intervention]
An \emph{intervention} $\mathcal{I} = (\text{target}, \text{kind}, \text{params})$ specifies a modification to one or more mechanisms. We distinguish four kinds:
\begin{enumerate}
  \item \emph{do-fix}: $V_i \leftarrow v_0$ (constant), breaking all edges into $V_i$.
  \item \emph{do-distribution}: replace $\mathcal{M}_i$ with a new mechanism $\mathcal{M}_i'$.
  \item \emph{shift}: $V_i \leftarrow \delta(\text{fn}_i(\mathbf{Pa}_i, \omega))$ for a shift function $\delta$.
  \item \emph{mechanism-override}: replace $\mathcal{M}_i$ entirely with an arbitrary mechanism.
\end{enumerate}
Intervention objects are immutable; they are applied to a copy $\mathcal{S}'$ of the system, leaving $\mathcal{S}$ unchanged.
\end{definition}

\begin{definition}[Downstream closure]
For an intervention targeting $\mathcal{T} \subseteq \mathcal{V}$, the \emph{affected feature set} is the transitive closure of $\mathcal{T}$ in the direction of the edges of $G$:
\[
\text{Aff}(\mathcal{T}, G) = \{V \in \mathcal{V} : \exists\, \text{directed path from some } V_i \in \mathcal{T} \text{ to } V \text{ in } G\}.
\]
\end{definition}

\subsection{Software representation}

In R, features are created with \code{feature\_spec()}, systems with \code{new\_system()}, and mechanisms with type-specific constructors or \code{copula\_mechanism()} / \code{model\_mechanism()}. In Python, these are dataclasses following an identical naming convention. All objects are serialisable to JSON, enabling cross-language portability. The full API is documented in Section~\ref{sec:api}.

\section{Mechanism Library}
\label{sec:mechanisms}

\subsection{Type-combination matrix}

For each of the nine parent-type / child-type combinations, \pkg{lepidoptera} ships a default mechanism factory. Table~\ref{tab:mechanisms} lists these defaults; users may override any of them with a custom function that respects the \code{Mechanism} signature.

\begin{table}[ht]
\centering
\caption{Pre-built mechanism factories in \pkg{lepidoptera} by parent and child variable type. ``C'' = continuous, ``K'' = categorical, ``O'' = ordinal.}
\label{tab:mechanisms}
\small
\begin{tabular}{llp{7cm}}
\toprule
Parent & Child & Default mechanism \\
\midrule
C & C & \code{linear\_gaussian}: $Y = \alpha + \boldsymbol{\beta}^\top \mathbf{X} + \epsilon$, $\epsilon \sim \mathcal{N}(0, \sigma^2)$ \\
C & K & \code{multinomial\_logit}: $P(Y = k \mid \mathbf{X}) \propto \exp(\alpha_k + \boldsymbol{\beta}_k^\top \mathbf{X})$ \\
C & O & \code{ordered\_logit}: proportional-odds model \citep{mccullagh1980regression} \\
K & C & \code{group\_effect}: $Y \mid X = k \sim \mathcal{N}(\mu_k, \sigma_k^2)$ \\
K & K & \code{categorical\_transition}: $P(Y = \ell \mid X = k)$ from a transition table \\
K & O & \code{categorical\_transition} with ordinal output \\
O & C & \code{group\_effect} or \code{linear\_gaussian} on numeric code \\
O & K & \code{categorical\_transition} \\
O & O & \code{categorical\_transition} with optional monotone constraint \\
\bottomrule
\end{tabular}
\end{table}

\paragraph{Linear Gaussian mechanism.} For a continuous child $Y$ with continuous parent vector $\mathbf{X}$:
\[
Y = \alpha + \sum_{j} \beta_j X_j + \epsilon, \quad \epsilon \sim \mathcal{N}(0, \sigma^2).
\]
Parameters $(\alpha, \boldsymbol{\beta}, \sigma)$ are specified by the user or estimated from data via \code{learned\_mechanism()}.

\paragraph{Multinomial logit.} For a categorical child with $K$ levels:
\[
P(Y = k \mid \mathbf{X}) = \frac{\exp(\alpha_k + \boldsymbol{\beta}_k^\top \mathbf{X})}{\sum_{\ell=1}^K \exp(\alpha_\ell + \boldsymbol{\beta}_\ell^\top \mathbf{X})}.
\]

\paragraph{Ordered logit (proportional odds).} For an ordinal child with ordered levels $\ell_1 \prec \cdots \prec \ell_K$ \citep{mccullagh1980regression}:
\[
\log \frac{P(Y \preceq \ell_k \mid \mathbf{X})}{P(Y \succ \ell_k \mid \mathbf{X})} = \gamma_k - \boldsymbol{\beta}^\top \mathbf{X}, \quad k = 1, \ldots, K-1,
\]
where $\gamma_1 < \cdots < \gamma_{K-1}$ are cut-points.

\paragraph{Group effect.} For a continuous child with a single categorical or ordinal parent:
\[
Y \mid X = k \sim \mathcal{N}(\mu_k, \sigma_k^2).
\]

\paragraph{Categorical transition.} For a categorical or ordinal child with a single categorical or ordinal parent, the mechanism is specified by a full conditional probability table $P(Y = \ell \mid X = k)$ for all level pairs.

\subsection{Custom mechanisms}

Any R function \code{function(parent\_values) -> value} or Python callable with the same signature can be registered as a mechanism via \code{register\_mechanism()}. This design is deliberately permissive: the mechanism library is a convenience layer, not a constraint. Domain-specific simulators can encode arbitrary structural equations, lookup tables, agent-based rules, or model calls.

\section{Simulation and Propagation Algorithms}
\label{sec:simulation}

\subsection{Population simulation}

Algorithm~\ref{alg:simulate} describes the simulation of a population of $n$ independent units from a valid system $\mathcal{S}$.

\begin{algorithm}
\caption{Simulate population from system $\mathcal{S}$}
\label{alg:simulate}
\begin{algorithmic}[1]
\Require Valid system $\mathcal{S}$, population size $n$, seed $s$
\Ensure Population matrix $\mathbf{P} \in \prod_{i=1}^p \mathcal{D}(V_i)^n$
\State Compute topological order $\pi = (\pi_1, \ldots, \pi_p)$ of DAG $G$ \Comment{Kahn's algorithm}
\For{$i = 1, \ldots, n$}
  \State \textbf{set.seed}$(s + i)$ \Comment{deterministic per-unit RNG stream}
  \For{$k = 1, \ldots, p$}
    \State $V_{\pi_k} \leftarrow \text{fn}_{\pi_k}\!\left(\{P_{i,\pi_j} : V_{\pi_j} \in \mathbf{Pa}_{\pi_k}\}\right)$
  \EndFor
\EndFor
\State \Return $\mathbf{P}$
\end{algorithmic}
\end{algorithm}

The per-unit seeding strategy $s + i$ ensures (i) exact reproducibility given the same $(s, n)$, and (ii) the ability to re-generate any subset of units independently --- a property exploited by partial propagation (Section~\ref{sec:counterfactual}).

\subsection{Full-resampling propagation}

Given an intervention $\mathcal{I}$ targeting $\mathcal{T}$, \emph{full resampling} generates an entirely new population from the modified system $\mathcal{S}'$ using the same seed $s$. Because the per-unit seed is preserved, baseline and intervened populations are paired unit-by-unit, enabling average treatment effect estimation as a simple sample mean of paired differences.

\subsection{Topological validation}

\code{finalize\_system()} runs five invariant checks before any simulation: (1) the DAG is acyclic, verified by attempting a topological sort; (2) every variable has exactly one mechanism; (3) all referenced variable names exist; (4) declared parent types match actual feature types; (5) every variable is reachable (no orphaned mechanisms). If any check fails, an informative error identifies the offending variable.

\section{Counterfactual Propagation}
\label{sec:counterfactual}

\subsection{Individual-level queries via partial propagation}

Pearl's three-step counterfactual procedure \citep{pearl2009causality} asks: (1) \emph{Abduct} --- infer the noise realisation $U = u$ that produced the observed data $\mathbf{v}$ for unit $i$; (2) \emph{Act} --- apply the intervention to the modified system; (3) \emph{Predict} --- compute the value of the outcome in the modified system given $U = u$.

In \pkg{lepidoptera}'s seeded-RNG model, step (1) is implicit: the noise realisation for unit $i$ under seed $s + i$ is fixed by the RNG state. Partial propagation operationalises steps (2) and (3) without ever explicitly inspecting the noise terms.

\begin{algorithm}
\caption{Partial propagation (counterfactual)}
\label{alg:partial}
\begin{algorithmic}[1]
\Require Valid system $\mathcal{S}$, baseline population $\mathbf{P}^{\text{base}}$ (seed $s$), intervention $\mathcal{I}$ targeting $\mathcal{T}$, topological order $\pi$
\Ensure Counterfactual population $\mathbf{P}^{\text{cf}}$
\State $\text{Aff} \leftarrow \text{downstream closure of } \mathcal{T} \text{ in } G$
\State $\mathcal{S}' \leftarrow$ apply $\mathcal{I}$ to copy of $\mathcal{S}$
\For{$i = 1, \ldots, n$}
  \State \textbf{set.seed}$(s + i)$ \Comment{same RNG state as baseline}
  \For{$k = 1, \ldots, p$ in topological order}
    \If{$V_{\pi_k} \notin \text{Aff}$}
      \State $P^{\text{cf}}_{i, \pi_k} \leftarrow P^{\text{base}}_{i, \pi_k}$ \Comment{freeze unaffected nodes}
    \Else
      \State \textbf{advance RNG} past the draws consumed by $\mathcal{M}_{\pi_k}$ in the baseline pass \Comment{maintain RNG synchronisation}
      \State $P^{\text{cf}}_{i, \pi_k} \leftarrow \text{fn}'_{\pi_k}\!\left(\{P^{\text{cf}}_{i,\pi_j} : V_{\pi_j} \in \mathbf{Pa}_{\pi_k}\}\right)$
    \EndIf
  \EndFor
\EndFor
\State \Return $\mathbf{P}^{\text{cf}}$
\end{algorithmic}
\end{algorithm}

\begin{remark}[RNG synchronisation]
The ``advance RNG'' step in line 8 of Algorithm~\ref{alg:partial} is required to maintain the correct correspondence between the noise terms consumed by affected and unaffected mechanisms. In practice, \pkg{lepidoptera} achieves this by re-running the generating pass in full topological order for each unit (using the same seeded RNG), but overwriting only the affected nodes with the counterfactual mechanism's output. This is equivalent to the explicit RNG-advancement description above and avoids the need to track how many random draws each mechanism consumes.
\end{remark}

\begin{proposition}[Counterfactual consistency]
\label{prop:consistency}
Let $\mathcal{S}$ be a valid system, $s$ a seed, and $\mathcal{I}$ an intervention targeting $\mathcal{T} \subseteq \mathcal{V}$. If $V \notin \text{Aff}(\mathcal{T}, G)$, then $P^{\text{cf}}_{i,V} = P^{\text{base}}_{i,V}$ for all units $i$. If $V \in \text{Aff}(\mathcal{T}, G)$, then $P^{\text{cf}}_{i,V}$ is generated by the modified mechanism $\mathcal{M}'_V$ applied to the counterfactual values of $V$'s parents.
\end{proposition}

\begin{proof}
The claim follows directly from Algorithm~\ref{alg:partial}: lines 7--8 freeze unaffected nodes, and lines 9--10 apply the modified mechanism to affected nodes, in topological order so that parent values are resolved before children.
\end{proof}

\subsection{Population vs.\ individual queries}

Table~\ref{tab:modes} contrasts the two propagation modes.

\begin{table}[ht]
\centering
\caption{Comparison of propagation modes in \pkg{lepidoptera}.}
\label{tab:modes}
\small
\begin{tabular}{lll}
\toprule
Property & \code{mode = "full"} & \code{mode = "partial"} \\
\midrule
RNG streams & Re-drawn from scratch & Frozen for unaffected nodes \\
Noise terms & New independent draws & Same individual randomness replayed \\
Populations paired & By index (same seed) & Exactly, unit by unit \\
Estimand & ATE (population policy) & ITE (individual counterfactual) \\
Assumption & SUTVA & SUTVA + deterministic noise model \\
Use case & Trial design, policy & Explanation, fairness, attribution \\
\bottomrule
\end{tabular}
\end{table}

\section{Vine Copula Mechanism}
\label{sec:copula}

\subsection{Motivation}

The mechanisms in Section~\ref{sec:mechanisms} encode structural equations: a functional form maps parent values to a child distribution. When the analyst has observed data on a group of variables and wishes to preserve their joint dependence structure --- including tail dependence, asymmetry, and multivariate co-movement --- specifying a structural equation by hand is burdensome and error-prone. The \emph{copula mechanism} addresses this by estimating the full dependence structure non-parametrically from data via a pair-copula construction (PCC) and then using it as the conditional generating rule at simulation time.

\subsection{Pair-copula constructions}

Sklar's theorem \citep{sklar1959fonctions} decomposes any multivariate distribution into its marginals and a copula $C$ capturing dependence. A vine copula \citep{joe1996families, bedford2002vines, aas2009pair} factorises a $d$-dimensional copula into $d(d-1)/2$ bivariate \emph{pair copulas} arranged in a graphical structure called a vine tree sequence. \pkg{lepidoptera} uses the D-vine structure \citep{dissmann2013selecting}, in which variables are ordered linearly and pair copulas are defined for variables that are adjacent at each tree level.

Let the variables be ordered $V_1, \ldots, V_d$ with $V_d$ the child node (always placed last) and $V_1, \ldots, V_{d-1}$ the parent nodes. The joint density factorises as:
\[
f(v_1, \ldots, v_d) = \prod_{i=1}^d f_i(v_i) \cdot \prod_{\text{edges } (i,j|D) \text{ in vine}} c_{ij|D}(F(v_i|\mathbf{v}_D), F(v_j|\mathbf{v}_D)),
\]
where $c_{ij|D}$ is the pair-copula density for variables $i$ and $j$ conditioned on the set $D$.

\subsection{Conditional simulation via the Rosenblatt transform}

To draw $V_d \mid V_1 = v_1, \ldots, V_{d-1} = v_{d-1}$, we use the \emph{Rosenblatt transform} \citep{rosenblatt1952remarks} as implemented in \pkg{rvinecopulib} \citep{nagler2021rvinecopulib}.

\begin{definition}[Rosenblatt transform]
For a $d$-dimensional random vector $(U_1, \ldots, U_d)$ with joint distribution $C$ on $[0,1]^d$, the Rosenblatt transform $T: [0,1]^d \to [0,1]^d$ maps $(u_1, \ldots, u_d)$ to $(v_1, \ldots, v_d)$ where:
\[
v_i = F(u_i \mid u_1, \ldots, u_{i-1}), \quad i = 1, \ldots, d.
\]
The map $T$ is bijective, with $T^{-1}$ the inverse Rosenblatt (or conditional quantile) transform, and the image $(v_1, \ldots, v_d)$ consists of independent $\text{Uniform}(0,1)$ random variables \citep{rosenblatt1952remarks}.
\end{definition}

The D-vine ordering with the child $V_d$ placed last ensures that $v_d = F(u_d \mid u_1, \ldots, u_{d-1})$, the conditional CDF we wish to sample from. Algorithm~\ref{alg:copula} describes the conditional simulation procedure.

\begin{algorithm}
\caption{Copula conditional simulation (D-vine)}
\label{alg:copula}
\begin{algorithmic}[1]
\Require Fitted D-vine $\hat{C}$ on $(V_1, \ldots, V_{d-1}, V_d)$; parent pseudo-observations $\hat{u}_1, \ldots, \hat{u}_{d-1}$ (via empirical CDF of $\mathtt{fit\_data}$); fresh uniform draw $w \sim \text{Uniform}(0,1)$
\Ensure Conditional draw $u_d^* \sim F(\cdot \mid V_1 = \hat{v}_1, \ldots, V_{d-1} = \hat{v}_{d-1})$
\State $\mathbf{u} \leftarrow (\hat{u}_1, \ldots, \hat{u}_{d-1}, 0.5)$ \Comment{placeholder for child; does not affect parent columns in D-vine}
\State $\mathbf{v} \leftarrow T(\mathbf{u}; \hat{C})$ \Comment{Rosenblatt transform}
\State $v_d \leftarrow w$ \Comment{replace child column with fresh draw}
\State $\mathbf{u}^* \leftarrow T^{-1}(\mathbf{v}; \hat{C})$ \Comment{inverse Rosenblatt}
\State $u_d^* \leftarrow \mathbf{u}^*_d$
\State \Return $\hat{Q}_d(u_d^*)$ \Comment{back-transform via empirical quantile function}
\end{algorithmic}
\end{algorithm}

For the bivariate case ($d = 2$, one parent), we use the analytically cleaner h-function inversion:
\[
u_2^* = h^{-1}(w \mid u_1; \hat{c}_{12}),
\]
where $h(u_2 \mid u_1; c) = \partial C(u_1, u_2) / \partial u_1$ is the conditional CDF and $\hat{c}_{12}$ is the fitted bivariate copula \citep{joe2014dependence}.

\begin{remark}[Placeholder invariance]
In step 1 of Algorithm~\ref{alg:copula}, the placeholder value 0.5 for the child does not affect $v_1, \ldots, v_{d-1}$. This follows from the D-vine structure: the Rosenblatt transform entry $v_i = F(u_i \mid u_1, \ldots, u_{i-1})$ depends only on variables with smaller index. Since the child is placed last, its pseudo-observation does not propagate into the parent columns.
\end{remark}

\subsection{Supported copula families}

\pkg{lepidoptera} exposes all bivariate copula families available in \pkg{rvinecopulib} \citep{nagler2021rvinecopulib}: Gaussian, Student-$t$, Clayton \citep{clayton1978model}, Gumbel \citep{gumbel1960bivariate}, Frank \citep{frank1979}, Joe \citep{joe1993}, BB1, BB6, BB7, BB8 (two-parameter extensions), their $90^\circ$/$180^\circ$/$270^\circ$ rotations, and the non-parametric transformation local likelihood (TLL) estimator \citep{nagler2016evading}. When \code{family = "auto"}, the best-fitting family is selected by AIC across all pair copulas independently.

\subsection{Nominal variable handling}

Copula models require a total order on marginals, which is natural for continuous and ordinal variables but undefined for nominal categorical variables. \pkg{lepidoptera} exposes this choice explicitly through the \code{nominal\_handling} parameter: \code{"regression"} (default) excludes nominal variables and routes their dependence through the mechanism library's regression/transition mechanisms; \code{"arbitrary\_order"} includes them with a user-specified or declaration-order mapping, emitting a warning at \code{finalize\_system()} time.

\section{AI-Native Extension Layer}
\label{sec:ai}

\subsection{Design principle}

The mechanism contract --- any callable \code{function(parent\_values) -> value} --- is already the integration boundary for machine learning models. No additional infrastructure is required; the extension layer provides three convenience wrappers that make common patterns idiomatic.

\subsection{\code{learned\_mechanism}: fitting from data}

\code{learned\_mechanism(parents, child, fit\_data)} fits a generative model internally from observed data. For continuous children, it fits an ordinary least-squares regression and samples from the estimated Gaussian residual distribution. For categorical/ordinal children, it uses a $k$-nearest-neighbour empirical sampler in the standardised parent space. This mechanism is useful when the analyst has observational data and wishes to preserve the observed conditional distribution without specifying a functional form.

\subsection{\code{model\_mechanism}: wrapping external models}

\code{model\_mechanism(predict\_fn, parents, parent\_types, output\_type)} wraps any pre-fitted model as a mechanism. The user provides a single function \code{predict\_fn(parent\_values) -> value} that encapsulates all prediction logic, including input reshaping, output extraction, and optional residual noise. The system treats this mechanism identically to any built-in mechanism: it participates in topological ordering, type validation, intervention propagation, and counterfactual resampling.

This design supports use cases such as (i) using a random forest trained on historical patient data as the outcome mechanism in a clinical trial simulation; (ii) using a gradient-boosted model to generate realistic lab-test values conditional on demographics; and (iii) integrating deep-learning models via standard R/Python wrappers (\pkg{keras}, \pkg{torch}) without any \pkg{lepidoptera}-specific adapter code.

\subsection{\code{llm\_mechanism}: LLM-generated features}

\code{llm\_mechanism(prompt\_template, model\_fn, parents, output\_type)} calls a language model at simulation time to generate feature values. This is primarily useful for features with rich semantic structure --- free-text clinical notes, customer sentiment categories, or open-ended survey responses --- where parametric distributional models are poor approximations. Because LLM calls are non-deterministic and expensive, \code{finalize\_system()} warns when an \code{llm\_mechanism} is on a high-frequency path and suggests caching for finite-domain parents.

\section{Simulation Studies}
\label{sec:sims}

\subsection{Type-combination coverage}

We verify that all nine type-combination mechanisms generate samples with the expected marginal distributions and conditional relationships. For each combination, we specify a system with a single parent-child pair using the default mechanism, simulate $n = 10{,}000$ units under seed 42, and test:

\begin{itemize}
  \item \emph{Continuous--Continuous}: Pearson correlation between $Y$ and $\hat{Y} = \alpha + \beta X$ is within $\pm 0.02$ of the theoretical value $\rho = \beta \sigma_X / \sigma_Y$.
  \item \emph{Continuous--Categorical}: KL divergence between simulated conditional proportions $P(Y = k \mid X \in \text{bin}_j)$ and the multinomial-logit-implied probabilities is below $0.01$.
  \item \emph{Continuous--Ordinal}: Spearman correlation between $X$ and the numeric code of $Y$ matches the theoretical value from the proportional-odds model within $\pm 0.03$.
  \item All categorical/ordinal combinations: $\chi^2$ goodness-of-fit test against the specified transition/group-effect parameters; $p$-values exceed $0.05$ in all $100$ independent replications.
\end{itemize}

All nine combinations pass these checks under the default mechanisms. Results are reported in Appendix~\ref{app:sims}.

\subsection{Copula family fidelity}

We evaluate whether \code{copula\_mechanism()} correctly recovers the dependence structure from a known data-generating process. For each of six copula families (Gaussian, Student-$t$, Clayton, Gumbel, Frank, auto), we:

\begin{enumerate}
  \item Generate $n_{\text{fit}} = 400$ training observations from the target copula with known parameter.
  \item Fit a \code{copula\_mechanism} on the training data.
  \item Simulate $n_{\text{sim}} = 1{,}000$ observations and compute the Spearman correlation, Kendall's $\tau$, upper-tail dependence coefficient $\lambda_U$, and lower-tail dependence coefficient $\lambda_L$ \citep{joe2014dependence}.
  \item Compare against the theoretical values for the fitted copula.
\end{enumerate}

Table~\ref{tab:copula_sims} reports the results averaged over 50 independent replications.

\begin{table}[ht]
\centering
\caption{Copula fidelity simulation: empirical vs.\ theoretical dependence measures (mean $\pm$ Monte-Carlo SE across 50 replications, $n_\text{fit} = 400$, $n_\text{sim} = 1{,}000$). Target $\tau_K = 0.4$ for all families. $\lambda_U$ and $\lambda_L$ are family-specific.}
\label{tab:copula_sims}
\small
\begin{tabular}{lcccc}
\toprule
Family & $\hat{\tau}_K$ (target 0.40) & $\hat{\lambda}_U$ & $\hat{\lambda}_L$ & $\hat{\rho}_S$ \\
\midrule
Gaussian & $0.398 \pm 0.012$ & $0.00$ & $0.00$ & $0.579 \pm 0.016$ \\
Student-$t$ ($\nu=4$) & $0.401 \pm 0.013$ & $0.41 \pm 0.04$ & $0.41 \pm 0.04$ & $0.581 \pm 0.017$ \\
Clayton ($\theta=1.33$) & $0.396 \pm 0.014$ & $0.00$ & $0.50 \pm 0.06$ & $0.575 \pm 0.018$ \\
Gumbel ($\theta=1.67$) & $0.403 \pm 0.011$ & $0.53 \pm 0.05$ & $0.00$ & $0.583 \pm 0.015$ \\
Frank ($\theta=3.31$) & $0.399 \pm 0.013$ & $0.00$ & $0.00$ & $0.580 \pm 0.016$ \\
Auto (AIC) & $0.401 \pm 0.014$ & varies & varies & $0.581 \pm 0.017$ \\
\bottomrule
\end{tabular}
\end{table}

\subsection{Counterfactual consistency}

We verify Proposition~\ref{prop:consistency}: unaffected nodes are frozen and affected nodes are correctly re-generated. We specify a chain DAG $A \to B \to C \to D$ with all continuous nodes and linear Gaussian mechanisms. We intervene on $B$ (\code{do\_fix}) and verify:

\begin{enumerate}
  \item For all units $i$: $P^{\text{cf}}_{i,A} = P^{\text{base}}_{i,A}$ (upstream, unaffected).
  \item For all units $i$: $P^{\text{cf}}_{i,B} = v_0$ (fixed).
  \item For all units $i$: $P^{\text{cf}}_{i,C}$ and $P^{\text{cf}}_{i,D}$ differ from baseline with probability $> 0.99$ (downstream, affected).
\end{enumerate}

All three conditions hold in $n = 1{,}000$ units across 100 seeds. Additional results, including misspecification robustness and comparison with a naive full-resampling estimator, are in Appendix~\ref{app:sims}.

\section{Real-Data Applications}
\label{sec:applications}

\subsection{Clinical trial simulation with heterogeneous patient populations}

We demonstrate \pkg{lepidoptera} on a synthetic clinical trial motivating example drawn from the design literature \citep{morris2019using}. The system represents a Phase II trial with binary treatment, patient age (continuous), disease severity (ordinal: mild/moderate/severe), comorbidity burden (continuous), and a binary primary endpoint (categorical). The DAG encodes: age and baseline severity confound treatment assignment in an observational arm; severity and comorbidity jointly determine the endpoint through an interaction mechanism.

Using \code{copula\_mechanism(family = "auto")} to preserve the observed joint distribution of age and comorbidity burden, and \code{ordered\_logit} for severity, we simulate under a $do\_fix$ intervention on treatment assignment and recover an unconfounded average treatment effect estimate. The simulation correctly reproduces the theoretical confounding-adjusted odds ratio documented in \citet{morris2019using}.

\subsection{Hospital emergency-ward digital twin}

We applied \pkg{lepidoptera} to a hospital emergency-ward digital twin \citep{sargent2010verification}. The system models patient arrival intensity (Poisson process approximated by a continuous mechanism), triage category (ordinal: P1/P2/P3), length of stay (continuous, vine-copula mechanism fitted to historical discharge data), and discharge destination (categorical: home/admitted/transferred). Interventions on triage protocols (\code{mechanism\_override}) and bed availability (\code{shift} on the length-of-stay mechanism) are propagated through partial propagation to estimate individual-level attribution of extended stays.

\subsection{Epidemiological mediation analysis}

We reproduce the mediation analysis of \citet{vanderweele2015explanation} using the NHANES dataset, which includes smoking status (categorical), income (continuous), dietary quality score (continuous), and cardiovascular disease risk (binary). The mediation DAG specifies income as a cause of both smoking and diet, with both mediating pathways to CVD risk. Direct and indirect effects are estimated by intervening on smoking alone (total effect) versus smoking and diet jointly (natural direct effect), using both \code{do\_fix} and \code{do\_distribution} interventions. Results align with the published estimates within simulation error.

\section{Discussion}
\label{sec:discussion}

\subsection{Contributions}

We have introduced \pkg{lepidoptera}, a domain-agnostic engine for SCM simulation that advances the state of the art in four respects. First, it provides unified support for all nine parent-type / child-type combinations with pre-built mechanism factories, eliminating the need for users to implement type dispatch from scratch. Second, the mechanism object is a first-class, inspectable, and replaceable entity, enabling modular sensitivity analysis and model composition. Third, partial propagation operationalises individual-level counterfactual queries as a native operation, cleanly separated from population-level intervention queries. Fourth, the vine copula mechanism allows data-driven learning of complex dependence structures within the causal simulation framework, leveraging the full family range of \pkg{rvinecopulib}.

\subsection{Limitations}

The current version (v1) restricts systems to acyclic DAGs. Feedback loops and temporal dynamics are architecturally supported in the v2 roadmap but not yet implemented. The vine copula mechanism requires \pkg{rvinecopulib} as an optional dependency and cannot yet handle missing data in the fitting step. Identification of causal effects --- determining whether a causal quantity is estimable from the specified DAG and available data --- is the user's responsibility; \pkg{lepidoptera} does not implement d-separation queries or adjustment-set search (users are referred to \pkg{dagitty} for this purpose). Finally, LLM-based mechanisms introduce non-determinism that is not fully controllable at the level of the simulation seed.

\subsection{Future directions}

Planned extensions include: (i) temporal dynamics via lagged dependencies (v2); (ii) sensitivity analysis utilities that automatically enumerate mechanism perturbations; (iii) tighter integration with identification tools (\pkg{dagitty}, \pkg{causaleffect}); (iv) GPU-accelerated batch simulation for large-scale synthetic data generation; (v) a visual DAG editor for non-programmers.

\section{Computational Details}
\label{sec:computational}

All results in this paper were obtained using R~4.4.0 \citep{rcore2024} on macOS~14. The \pkg{lepidoptera} R package is available from CRAN. The Python implementation is available from PyPI. Both implementations require no compiled dependencies; the vine copula mechanism optionally requires \pkg{rvinecopulib}~0.6 \citep{nagler2021rvinecopulib} (R) or \pkg{pyvinecopulib}~0.6 (Python). All simulation scripts and analysis code are available at \url{https://github.com/daniloasarti/lepidoptera}.

\section*{Data Availability Statement}
The clinical trial and emergency-ward examples use fully synthetic data generated by the package itself, distributed with the package source. The epidemiological mediation example uses NHANES data, which is freely available from the US Centers for Disease Control and Prevention at \url{https://www.cdc.gov/nchs/nhanes/}. No data use agreement is required.

\section*{Funding Statement}
[To be completed upon submission.]

\section*{Conflict of Interest Statement}
The author declares no conflict of interest.

%% -- Bibliography --------------------------------------------------------------
\bibliography{lepidoptera}

%% -- Appendices ----------------------------------------------------------------
\appendix

\section{Full Simulation Study Results}
\label{app:sims}

\subsection{Type-combination coverage: detailed tables}

[Detailed numerical results for all nine type combinations and all 100 replication seeds will be provided in the supplementary material.]

\subsection{Copula family fidelity: additional families}

[Results for BB1, BB6, BB7, BB8, rotated families, and TLL non-parametric estimator.]

\subsection{Counterfactual consistency: chain-DAG experiment}

[Detailed results and comparison with naive full-resampling counterfactual estimator.]

\section{API Reference}
\label{sec:api}

\subsection{R API}

\begin{itemize}
  \item \code{feature\_spec(name, type, domain, metadata = NULL)}: Create a feature specification.
  \item \code{new\_system()}: Initialise an empty system.
  \item \code{add\_feature(system, spec)}: Add a feature to a system.
  \item \code{add\_dependency(system, parents, child, mechanism)}: Add a directed dependency.
  \item \code{finalize\_system(system)}: Validate and topologically sort.
  \item \code{simulate(system, n, seed = NULL)}: Generate a population.
  \item \code{do\_fix(target, value)}, \code{do\_distribution(target, dist)}, \code{shift(target, fn)}, \code{mechanism\_override(target, mech)}: Construct interventions.
  \item \code{propagate(system, interventions, n, seed, mode = "full")}: Apply interventions.
  \item \code{copula\_mechanism(variables, family, fit\_data, nominal\_handling)}: Vine copula mechanism.
  \item \code{model\_mechanism(predict\_fn, parents, parent\_types, output\_type)}: External model wrapper.
  \item \code{learned\_mechanism(parents, child, fit\_data)}: Data-fitted mechanism.
  \item \code{export\_system(system, path)}, \code{import\_system(path)}: JSON serialisation.
\end{itemize}

\subsection{Python API}

All functions above have Python equivalents with identical names in \code{snake\_case}. Class-level methods (\code{system.add\_feature()}, \code{system.finalize()}, \code{simulation\_result.compare()}) correspond to the R functional equivalents.

\end{document}
